%! Introducing Vectors — Unit 1, Lesson 1.0 — Warm-Up
\documentclass[10pt]{article}
\usepackage{linalg-article}
\usepackage{linalg-boxes}
\newcommand{\TallMath}[1]{$\displaystyle #1\rule[-1.4em]{0pt}{3.2em}$}

\begin{document}

\pageheader{Unit 1, Lesson 1.0}{Warm-Up}
\namedateperiod

\vspace{0.12in}

\begin{notesbox}{Getting Ready: Skills We'll Use Today}
\small These three skills power everything in today's lesson. Answer quickly.

\vspace{4pt}
\textbf{1. Signed-number arithmetic.} Evaluate.
\begin{enumerate}[label=(\alph*), leftmargin=2.0em, itemsep=6pt]
  \item $3 + (-7) = \blank{2.2cm}$
  \item $-4 - (-9) = \blank{2.2cm}$
  \item $2(-5) = \blank{2.2cm}$
\end{enumerate}

\vspace{6pt}
\textbf{2. The Pythagorean theorem.} A right triangle has legs of length $6$ and $8$.
How long is the hypotenuse? Show the step.
\writelines{2}

\vspace{4pt}
\textbf{3. Reading the plane.} On the grid, point $P$ is marked.
\begin{center}
\begin{tikzpicture}[scale=0.5]
  \draw[step=1, linegray!60, thin] (-0.4,-0.4) grid (5.4,4.4);
  \draw[->, thick, charcoal] (0,0)--(5.6,0) node[right]{\scriptsize $x$};
  \draw[->, thick, charcoal] (0,0)--(0,4.6) node[above]{\scriptsize $y$};
  \foreach \x in {1,2,3,4,5} \node[below, font=\scriptsize] at (\x,-0.1){\x};
  \foreach \y in {1,2,3,4} \node[left, font=\scriptsize] at (-0.1,\y){\y};
  \fill[burgundy] (5,3) circle (3pt) node[above right]{\small $P$};
\end{tikzpicture}
\end{center}
What are the coordinates of $P$? \; $\big(\,\blank{1.0cm},\ \blank{1.0cm}\,\big)$
\end{notesbox}

\end{document}
